% GENERATED FILE. Edit canonical lessons, then run npm run generate:books.
% canonical-source-hash: fnv1a64:9c0a49e5608dbc2e
% canonical-lessons: RU-C170-dolzhnost, RU-C170-zarplata, RU-C170-nalog, RU-C170-schyot, RU-C170-nalichnyye

\chapter{Post, Salary, Tax, Bill, Cash}
\label{ch:ru-post-salary-tax-bill-cash}
\hlchaptermodality{170}

\begin{chapteropening}
\textbf{By the end of this chapter:} \emph{I can say a post, a salary, a tax, a bill and cash.}

Complete the last lesson of chapter 170: I can say a post, a salary, a tax, a bill and cash.
\end{chapteropening}

\section[\texorpdfstring{\ru{должность} (dólzhnost)}{должность (dólzhnost)}]{\ru{должность} (dólzhnost) — a post, a position}

\label{lesson:RU-C170-dolzhnost}



\begin{quote}
Before the new one: say the Russian for the truth, then the Russian for a lie.
\end{quote}

\subsection*{You'll want to know: \ru{должность}}
\textbf{\ru{должность}} (\emph{dólzhnost}) — \textquotedblleft{}a post, a position\textquotedblright{}.

\begin{cousinweb}
\textbf{\ru{должность}} comes from \emph{\ru{долг}}, \textquotedblleft{}a debt, a duty\textquotedblright{}, and first meant \textquotedblleft{}an obligation\textquotedblright{} (Vasmer). Vasmer relates \emph{\ru{долг}} to Gothic \emph{dulgs}, \textquotedblleft{}a debt\textquotedblright{}, and Old Irish \emph{dligid}, \textquotedblleft{}a duty\textquotedblright{}.
\end{cousinweb}

\begin{grammarlens}[title={What you've built}]
One more word for this chapter.
\end{grammarlens}

\subsection*{Your turn}
\begin{itemize}
\raggedright
  \item \emph{Say it:} \emph{dólzhnost}
  \item \emph{Say it:} \emph{dólzhnost}, once more
  \item \emph{Say it:} say \emph{dólzhnost}
  \item \emph{Recall:} say the Russian for the truth, then the Russian for a lie, then say \emph{dólzhnost} again
\end{itemize}

\begin{tcolorbox}[breakable,colback=teal!4,colframe=teal!35!black,title={Before you move on}]
What does \ru{должность} mean? (\textquotedblleft{}A post, a position\textquotedblright{}.) Say it once more.
\end{tcolorbox}
\section[\texorpdfstring{\ru{зарплата} (zarpláta)}{зарплата (zarpláta)}]{\ru{зарплата} (zarpláta) — a salary}

\label{lesson:RU-C170-zarplata}



\begin{quote}
Before the new one: say the Russian for a lie, then the Russian for a post.
\end{quote}

\subsection*{You'll want to know: \ru{зарплата}}
\textbf{\ru{зарплата}} (\emph{zarpláta}) — \textquotedblleft{}a salary\textquotedblright{}.

\begin{cousinweb}
\textbf{\ru{зарплата}} is \emph{\ru{заработная} \ru{плата}}, \textquotedblleft{}earned pay\textquotedblright{}, cut short. Vasmer explains \emph{\ru{плата}}, \textquotedblleft{}pay\textquotedblright{}, from \emph{\ru{плат}}, \textquotedblleft{}a cloth\textquotedblright{}, since pieces of cloth once served as money.
\end{cousinweb}

\begin{grammarlens}[title={What you've built}]
One more word for this chapter.
\end{grammarlens}

\subsection*{Your turn}
\begin{itemize}
\raggedright
  \item \emph{Say it:} \emph{zarpláta}
  \item \emph{Say it:} \emph{zarpláta}, once more
  \item \emph{Say it:} say \emph{zarpláta}
  \item \emph{Recall:} say the Russian for a lie, then the Russian for a post, then say \emph{zarpláta} again
\end{itemize}

\begin{tcolorbox}[breakable,colback=teal!4,colframe=teal!35!black,title={Before you move on}]
What does \ru{зарплата} mean? (\textquotedblleft{}A salary\textquotedblright{}.) Say it once more.
\end{tcolorbox}
\section[\texorpdfstring{\ru{налог} (nalóg)}{налог (nalóg)}]{\ru{налог} (nalóg) — a tax}

\label{lesson:RU-C170-nalog}



\begin{quote}
Before the new one: say the Russian for a post, then the Russian for a salary.
\end{quote}

\subsection*{You'll want to know: \ru{налог}}
\textbf{\ru{налог}} (\emph{nalóg}) — \textquotedblleft{}a tax\textquotedblright{}.

\begin{cousinweb}
\textbf{\ru{налог}} is \emph{\ru{на}-}, \textquotedblleft{}on\textquotedblright{}, with \emph{-\ru{лог}}, from the root of \emph{\ru{лежать}}, \textquotedblleft{}to lie\textquotedblright{}, and \emph{\ru{положить}}, \textquotedblleft{}to lay\textquotedblright{}: a tax is what is laid on you.
\end{cousinweb}

\begin{grammarlens}[title={What you've built}]
One more word for this chapter.
\end{grammarlens}

\subsection*{Your turn}
\begin{itemize}
\raggedright
  \item \emph{Say it:} \emph{nalóg}
  \item \emph{Say it:} \emph{nalóg}, once more
  \item \emph{Say it:} say \emph{nalóg}
  \item \emph{Recall:} say the Russian for a post, then the Russian for a salary, then say \emph{nalóg} again
\end{itemize}

\begin{tcolorbox}[breakable,colback=teal!4,colframe=teal!35!black,title={Before you move on}]
What does \ru{налог} mean? (\textquotedblleft{}A tax\textquotedblright{}.) Say it once more.
\end{tcolorbox}
\section[\texorpdfstring{\ru{счёт} (schyot)}{счёт (schyot)}]{\ru{счёт} (schyot) — a bill; an account}

\label{lesson:RU-C170-schyot}



\begin{quote}
Before the new one: say the Russian for a salary, then the Russian for a tax.
\end{quote}

\subsection*{You'll want to know: \ru{счёт}}
\textbf{\ru{счёт}} (\emph{schyot}) — \textquotedblleft{}a bill; an account\textquotedblright{}.

\begin{cousinweb}
\textbf{\ru{счёт}} is \emph{\ru{с}-} with \emph{\ru{чёт}}, \textquotedblleft{}a count\textquotedblright{}, from \emph{\ru{считать}}, \textquotedblleft{}to count\textquotedblright{}, on the root of \emph{\ru{читать}}, \textquotedblleft{}to read\textquotedblright{}, which once meant \textquotedblleft{}to count\textquotedblright{}.
\end{cousinweb}

\begin{grammarlens}[title={What you've built}]
One more word for this chapter.
\end{grammarlens}

\subsection*{Your turn}
\begin{itemize}
\raggedright
  \item \emph{Say it:} \emph{schyot}
  \item \emph{Say it:} \emph{schyot}, once more
  \item \emph{Say it:} say \emph{schyot}
  \item \emph{Recall:} say the Russian for a salary, then the Russian for a tax, then say \emph{schyot} again
\end{itemize}

\begin{tcolorbox}[breakable,colback=teal!4,colframe=teal!35!black,title={Before you move on}]
What does \ru{счёт} mean? (\textquotedblleft{}A bill; an account\textquotedblright{}.) Say it once more.
\end{tcolorbox}
\section[\texorpdfstring{\ru{наличные} (nalíchnyye)}{наличные (nalíchnyye)}]{\ru{наличные} (nalíchnyye) — cash}

\label{lesson:RU-C170-nalichnyye}



\begin{quote}
Before the new one: say the Russian for a tax, then the Russian for a bill; an account.
\end{quote}

\subsection*{You'll want to know: \ru{наличные}}
\textbf{\ru{наличные}} (\emph{nalíchnyye}) — \textquotedblleft{}cash\textquotedblright{}.

\begin{cousinweb}
\textbf{\ru{наличные}} is short for \emph{\ru{наличные} \ru{деньги}}, \textquotedblleft{}money at hand\textquotedblright{}, from \emph{\ru{налицо}}, \textquotedblleft{}present, in view\textquotedblright{}, \emph{\ru{на}}, \textquotedblleft{}on\textquotedblright{}, with \emph{\ru{лицо}}, \textquotedblleft{}face\textquotedblright{}.
\end{cousinweb}

\begin{grammarlens}[title={What you've built}]
One more word for this chapter.
\end{grammarlens}

\subsection*{Your turn}
\begin{itemize}
\raggedright
  \item \emph{Say it:} \emph{nalíchnyye}
  \item \emph{Say it:} \emph{nalíchnyye}, once more
  \item \emph{Say it:} say \emph{nalíchnyye}
  \item \emph{Recall:} say the Russian for a tax, then the Russian for a bill; an account, then say \emph{nalíchnyye} again
\end{itemize}

\begin{tcolorbox}[breakable,colback=teal!4,colframe=teal!35!black,title={Before you move on}]
What does \ru{наличные} mean? (\textquotedblleft{}Cash\textquotedblright{}.) Say it once more.
\end{tcolorbox}
